\documentclass[11pt]{article}
\usepackage[margin=1in]{geometry}
\usepackage{amsmath,amssymb,amsthm}
\usepackage[hidelinks]{hyperref}

\newtheorem{theorem}{Theorem}
\newtheorem{lemma}{Lemma}
\newtheorem{corollary}{Corollary}
\newtheorem*{theorem*}{Theorem}

\title{A Polynomial-Space $O(1.3069^n)$ Algorithm for\\ 3-Path Vertex Cover on Subcubic Graphs}
\author{}
\date{}

\begin{document}
\maketitle

\begin{abstract}
A known lower bound on the dissociation number gives
$\operatorname{diss}(G) \ge n/\lceil(\Delta+1)/2\rceil$ for every graph of maximum
degree $\Delta$; for subcubic graphs this is $\operatorname{diss}(G) \ge n/2$.
Hence the minimum 3-path vertex cover has size at most $n/2$. Combining this bound
with the polynomial-space parameterized algorithm of \v{C}erven\'y and Such\'y
(arXiv:2111.05896), which runs in $O^*(1.708^k)$ time, gives an
$O^*(1.708^{n/2}) = O(1.3069^n)$-time polynomial-space algorithm for 3-path vertex
cover on subcubic graphs. This improves both the previous polynomial-space bound
$O(1.4613^n)$ of Chang et al.\ (2018) and the exponential-space bound
$O(1.3659^n)$ of Xiao and Kou (2017). Conditionally on the 2026 preprint of Gaspers
(arXiv:2611.05146), we also record $O(1.2823^n)$ for all subcubic graphs and
$O(1.2812^n)$ for triangle-free subcubic graphs, both in polynomial space.
\end{abstract}

\section{Definitions and known results}

A \emph{3-path} is a path on three vertices. A \emph{3-path vertex cover} is a set
$C \subseteq V$ such that $G-C$ contains no 3-path, equivalently $G-C$ has maximum
degree at most $1$. A \emph{dissociation set} is a set $S \subseteq V$ with
$\Delta(G[S]) \le 1$. The two notions are complementary: $S$ is a dissociation set if
and only if $V \setminus S$ is a 3-path vertex cover. Let $k^*(G)$ be the minimum size
of a 3-path vertex cover of $G$.

Known exact algorithms for an $n$-vertex graph are
$O^*(1.5171^n)$ (Kardo\v{s}, Katreni\v{c}, and Schiermeyer),
$O^*(1.4658^n)$ (Chang et al.),
$O^*(1.4613^n)$ time and polynomial space (Chang et al., 2018),
$O^*(1.3659^n)$ time with exponential space (Xiao and Kou, 2017).

For the parameterized version, the running time was improved successively to
$O^*(2^k)$ (Tu), $O^*(1.8127^k)$ (Katreni\v{c}), $O^*(1.7485^k)$ (Xiao and Kou, polynomial
space), $O^*(1.713^k)$ (Tsur), and $O^*(1.708^k)$ (\v{C}erven\'y and Such\'y, 2021,
branch-and-reduce, polynomial space).

\section{A dissociation set of size at least $n/2$}

We use the following known lower bound on the dissociation number; it is due to
Bre\v{s}ar, Kardo\v{s}, Katreni\v{c}, and Semani\v{s}in (2011) (Theorem~6 and
Corollary~7), and is also restated in Bock, Pardey, Penso, and Rautenbach
(J.\ Graph Theory, 2023).

\begin{lemma}
If $G$ has maximum degree $\Delta$, then
$\operatorname{diss}(G) \ge n/\lceil(\Delta+1)/2\rceil$.
\end{lemma}

\begin{proof}
By Theorem~6 of Bre\v{s}ar et al., the vertex set of a graph of maximum degree
$\Delta$ can be partitioned, in $O(|E(G)|)$ time, into $\lceil(\Delta+1)/2\rceil$
sets each inducing a subgraph of maximum degree at most $1$. The largest set has
size at least $n/\lceil(\Delta+1)/2\rceil$ and is a dissociation set.
\end{proof}

\begin{corollary}
Every subcubic graph satisfies $\operatorname{diss}(G) \ge n/2$, and hence
$k^*(G) \le n/2$.
\end{corollary}

\begin{proof}
For $\Delta \le 3$ we have $\lceil(\Delta+1)/2\rceil \le 2$, so
$\operatorname{diss}(G) \ge n/2$. The complement of a maximum dissociation set is a
3-path vertex cover of size $n-\operatorname{diss}(G) \le n/2$.
\end{proof}

\section{The algorithm}

\begin{enumerate}
\item Let $k_0 = \lfloor n/2 \rfloor$. By Corollary 2, $k^*(G) \le k_0$.
\item For $k = 0,1,\dots,k_0$, run the polynomial-space decision algorithm
      \texttt{P3VC}$(G,k)$ of \v{C}erven\'y and Such\'y.
\item Return the smallest $k$ for which \texttt{P3VC}$(G,k)$ answers yes.
\end{enumerate}

\textbf{Correctness.} \texttt{P3VC} is exact, and $k^*(G) \le k_0$; the loop therefore
returns $k^*(G)$.

\textbf{Time.} The $k$-th call takes $O^*(1.708^k)$ time. The sum over
$k=0,\dots,k_0$ is dominated by the last term, hence
\[
  O^*(1.708^{k_0}) \;\le\; O^*(1.708^{n/2}) \;=\; O(1.3069^n).
\]
($1.708^{1/2} = 1.30691\ldots$.)

\textbf{Space.} The algorithm of \v{C}erven\'y and Such\'y is a branch-and-reduce
algorithm and uses polynomial space; the loop only stores the current instance. Hence
the whole algorithm uses polynomial space.

\section{Comparison}

\begin{center}
\begin{tabular}{lll}
\hline
Algorithm & Time & Space \\
\hline
Kardo\v{s} et al.\ 2011 & $O^*(1.5171^n)$ & polynomial \\
Chang et al. & $O^*(1.4658^n)$ & polynomial \\
Chang et al.\ 2018 & $O^*(1.4613^n)$ & polynomial \\
Xiao--Kou 2017 & $O^*(1.3659^n)$ & exponential \\
\textbf{This note} & $\boldsymbol{O(1.3069^n)}$ & \textbf{polynomial} \\
\hline
\end{tabular}
\end{center}

The new bound improves both the polynomial-space bound $O(1.4613^n)$ and the
exponential-space bound $O(1.3659^n)$.

\section{Conditional improvements via Max 2-CSP}

The bounds in this section rely on the following 2026 preprint, which has not yet
been peer-reviewed. We state them conditionally; the \(O(1.3069^n)\) bound of the
previous sections remains unconditional.

\begin{theorem*}[Gaspers, arXiv:2611.05146, Theorem~C.3(a)]
For every \(r\ge2\), cubic Max \((r,2)\)-CSP can be solved in polynomial space and
time \(r^{\gamma N+o(N)}\), where
\[
  \gamma=\frac{2c_3}{1+4c_3},\qquad
  c_3=\frac{3}{4\pi}\arccos\frac56<0.1398,
\]
so \(\gamma<0.179341\).
\end{theorem*}

\subsection*{A four-state encoding}

Let \(G\) be subcubic with no isolated vertices. For each \(v\in V(G)\) create a
variable \(x_v\in\{\bot\}\cup N(v)\); since \(\deg(v)\le3\), the domain has size at
most \(4\). The unary score is \(w_v(\bot)=0\) and \(w_v(u)=1\) for \(u\in N(v)\).
For each edge \(uv\) the binary score is \(0\) if \(x_u=\bot\) or \(x_v=\bot\) or
\((x_u,x_v)=(v,u)\), and \(-(n+1)\) otherwise. This is a Max 2-CSP whose constraint
graph is \(G\).

\begin{lemma}
The maximum score of the encoding equals \(\operatorname{diss}(G)\).
\end{lemma}

\begin{proof}
If \(S\) is a dissociation set, set \(x_v\) to the unique neighbour of \(v\) in \(S\)
when it exists, to an arbitrary neighbour when \(v\in S\) has no neighbour in \(S\),
and to \(\bot\) when \(v\notin S\). Every edge with both endpoints in \(S\) joins two
paired vertices, so no penalty is paid and the score is \(|S|\). Conversely, if an
assignment pays no penalty, let \(S=\{v:x_v\ne\bot\}\). If \(v\in S\) had two
neighbours in \(S\), then \(x_v\) would have to equal both of them, impossible; hence
\(S\) is a dissociation set and the score is \(|S|\). Finally, any assignment that
pays a penalty scores at most \(n-(n+1)<0\), while the all-\(\bot\) assignment
scores \(0\); so an optimum pays no penalty. Isolated vertices are put into \(S\)
separately and removed.
\end{proof}

Eliminating a variable of degree at most \(2\) is exact: for each assignment of its
neighbours, take the maximum over its at most four values and add the resulting
potential to the neighbours. The constraint graph stays subcubic, and after all
degree-\(0,1,2\) variables are eliminated the remaining graph is cubic on
\(N\le n_3\) variables, where \(n_3\) is the number of degree-\(3\) vertices of
\(G\). Recording the argmax in each elimination recovers an optimal assignment in
polynomial time and space.

\begin{theorem*}[conditional, all subcubic graphs]
Assume the Gaspers theorem above. Then 3-path vertex cover on subcubic graphs can be
solved in polynomial space and time \(O(1.2823^n)\).
\end{theorem*}

\begin{proof}
Encode \(G\) as above, eliminate degree-\(\le2\) variables, and run the algorithm of
the Gaspers theorem on the cubic core with \(r=4\). Since
\(4^\gamma=1.282254\ldots<1.2823\) and \(N\le n_3\le n\), the running time is
\(O(1.2823^n)\). The solution is recovered by self-reduction: fix variables one at a
time and rerun the optimisation, at most four times per variable. This adds a
polynomial factor and keeps polynomial space.
\end{proof}

\subsection*{Local branching for triangle-free graphs}

Throughout this subsection \(G\) is triangle-free and subcubic. Let \(n_3\) be the
number of degree-\(3\) vertices and \(\mu(G)=0.95n_3+0.05n\). A vertex \(v\) is a
\emph{branching point} if every vertex within distance at most \(3\) of \(v\) has
degree \(3\). Let \(B_2(v)\) be the set of vertices within distance at most \(2\) of
\(v\).

For a branching point \(v\) we branch as follows: delete \(v\); keep \(v\) isolated
(delete \(N[v]\) and add \(1\)); or pair \(v\) with one of its three neighbours \(u\)
(delete \(N[v]\cup N[u]\) and add \(2\)). The following table gives the guaranteed
decrease of \(n_3\) and of \(\mu\) in the five branches; the branch deleting \(v\) is
first, the branch keeping \(v\) isolated second, and the three pairing branches last.

\begin{center}
\small
\begin{tabular}{llll}
\hline
Case & \(n_3\) decrease & \(\mu\) decrease & \(\sum r^{-\Delta\mu}\) at \(r=1.2812\) \\
\hline
\(|B_2(v)|\ge8\) & \(4,8,8,8,8\) & \(3.85,7.8,7.9,7.9,7.9\) & \(0.9535\) \\
\(|B_2(v)|\le7\), no equal neighbourhoods & \(4,7,8,8,8\) & \(3.85,6.85,7.9,7.9,7.9\) & \(0.991939\) \\
equal neighbourhoods & two branches, \(n_3\) decrease \(\ge5\) & --- & \(<0.591\) \\
\hline
\end{tabular}
\end{center}

In the second line, the three neighbours of \(v\) have second-layer neighbours
forming three distinct pairs \(\{a,b\},\{b,c\},\{a,c\}\); triangle-freeness makes
\(a,b,c\) pairwise non-adjacent, and every pairing branch removes the seven
degree-\(3\) vertices of the ball and at least one further degree-\(3\) vertex. In the
third line, if two neighbours of \(v\) have the same neighbourhood, some optimum
keeps both or deletes both; the two branches each decrease \(n_3\) by at least \(5\)
and their \(r^{-\Delta\mu}\) sum is below \(0.591\).

If no branching point exists, every degree-\(3\) vertex is within distance \(3\) of a
low-degree vertex. Such a vertex covers at most \(2+4+8=14\) degree-\(3\) vertices, so
\(n_3\le14(n-n_3)\), and the backend of the conditional all-subcubic bound costs at
most
\[
  1.2823^{n_3}\le1.2823^{14\mu/14.05}
  =1.2811658\ldots^{\mu}<1.2812^{\mu}.
\]
At a branching point the five-branch inequality is below \(1\) at \(r=1.2812\)
(\(0.9535\) and \(0.991939\) in the two cases), and the two-branch inequality is
below \(1\). By induction on \(\mu\), the total running time is
\(O(1.2812^{\mu(G)})=O(1.2812^n)\). A depth-first traversal records the best branch
at each node and recovers an optimal solution in polynomial space.

\begin{theorem*}[conditional, triangle-free subcubic graphs]
Assume the Gaspers theorem above. Then 3-path vertex cover on triangle-free subcubic
graphs can be solved in polynomial space and time \(O(1.2812^n)\).
\end{theorem*}

\subsection*{Diamonds and \(K_4\)}

\begin{lemma}
Let \(D\) induce a diamond \(K_4-e\) whose two degree-\(3\) vertices have no
neighbours outside \(D\), or let \(D\) be a \(K_4\) component. Then
\(\psi(G)=2+\psi(G-D)\).
\end{lemma}

\begin{proof}[Proof sketch]
Any dissociation set keeps at most two vertices of \(D\). Conversely, an exchange
argument shows that an optimum can be chosen to keep the two adjacent degree-\(3\)
vertices of the diamond (or any two vertices of the \(K_4\)); deleting the other two
loses nothing.
\end{proof}

Repeating these reductions leaves a graph that is triangle-free or whose triangles
are vertex-disjoint; the bound \(O(1.2812^n)\) applies to the triangle-free case,
while the general subcubic bound remains \(O(1.2823^n)\).

\smallskip
\noindent\textbf{Remark.}
The case analysis of this section is a proof sketch: the numerical inequalities have
been checked, and the encoding and the diamond/\(K_4\) reduction have been verified by
exhaustive small-case tests, but the structural case analysis has not been
formalised. The bounds are conditional on the Gaspers preprint. The bound
\(O(1.2812^n)\) is specific to triangle-free (or diamond/\(K_4\)-reducible) graphs; it
is not claimed for all subcubic graphs.

\section{Remarks and limitations}

\begin{itemize}
\item The result is specific to graphs of maximum degree 3.
\item The improvement over the exponential-space bound is about $4.3\%$.
\item The lower bound $\operatorname{diss}(G) \ge n/2$ is tight: $K_4$ has
      $\operatorname{diss}(K_4)=2=n/2$.
\item The argument is black-box with respect to the parameterized algorithm. If the
      parameterized bound $1.708^k$ is improved, the same construction improves
      accordingly: a parameterized bound $c^k$ and the bound $k \le n/2$ give
      $c^{n/2}$.
\end{itemize}

\section{Novelty check}

A web search for the bound $1.3069$ and related queries did not find a prior
$1.3069^n$ bound. The ingredients are known: the lower bound
$\operatorname{diss}(G) \ge n/\lceil(\Delta+1)/2\rceil$ (see Bre\v{s}ar et al., 2011;
Bock et al., 2023) and
the parameterized algorithm of \v{C}erven\'y and Such\'y ($O^*(1.708^k)$). The hybrid
is elementary and may be folklore; novelty should be verified against the latest
literature before submission.

\begin{thebibliography}{9}

\bibitem{gaspers26}
S.~Gaspers.
Path decompositions of sparse graphs via Gaussian sweeps and faster exponential
algorithms.
arXiv:2611.05146, 2026.

\bibitem{bkks11}
B.~Bre\v{s}ar, F.~Kardo\v{s}, J.~Katreni\v{c}, and G.~Semani\v{s}in.
Minimum $k$-path vertex cover.
\emph{Discrete Applied Mathematics} 159(12) (2011), 1189--1195.

\bibitem{bppr23}
F.~Bock, J.~Pardey, L.~D.~Penso, and D.~Rautenbach.
A bound on the dissociation number.
\emph{Journal of Graph Theory} 103(4) (2023), 661--673.

\bibitem{cs23}
R.~\v{C}erven\'y and O.~Such\'y.
Generating faster algorithms for $d$-Path Vertex Cover.
In \emph{WG 2023}, LNCS 14093, Springer, 2023, pp.~157--171.
(Full version: arXiv:2111.05896.)

\bibitem{tsur}
D.~Tsur.
Parameterized algorithm for 3-path vertex cover.
arXiv:1809.02636, 2018.

\bibitem{chang18}
M.-S.~Chang, L.-H.~Chen, L.-J.~Hung, Y.-Z.~Liu, P.~Rossmanith, and S.~Sikdar.
Moderately exponential time algorithms for the maximum bounded-degree-1 set problem.
\emph{Discrete Applied Mathematics}, 2018.

\bibitem{xk17}
M.~Xiao and S.~Kou.
Exact algorithms for the maximum dissociation set and minimum 3-path vertex cover
problems.
\emph{Theoretical Computer Science} 657 (2017), 86--97.

\bibitem{xk16}
M.~Xiao and S.~Kou.
Kernelization and Parameterized Algorithms for 3-Path Vertex Cover.
arXiv:1608.07022, 2016.

\bibitem{kks11}
F.~Kardo\v{s}, J.~Katreni\v{c}, and I.~Schiermeyer.
On computing the minimum 3-path vertex cover and dissociation number of graphs.
\emph{Theoretical Computer Science} 412 (2011), 7009--7017.

\end{thebibliography}

\end{document}
